% Source: https://artofproblemsolving.com/community/c3h1572701_hard_problem
% Stable identity: aops-topic-1572701-q01
Consider a collection with the largest product.
\begin{itemize}
\item No part is $1$ (unless the whole sum is $1$): merging a $1$ with another part $k$ gives $k+1>1\cdot k$.
\item No part is $5$ or more: a part $k\ge5$ can be split into $2$ and $k-2$, and $2(k-2)=k+(k-4)>k$.
\item There are at most two $2$s (and a $4$ can be replaced by $2+2$ with the same product): three $2$s can be replaced by two $3$s, since $3\cdot3=9>8=2\cdot2\cdot2$.
\end{itemize}
So, writing every $4$ as $2+2$, an optimal collection consists of $3$s and at most two $2$s. Since $19=3\cdot5+4$, and the number of $2$s is determined by $19\bmod3=1$ (one $2$ would leave $17$, not a multiple of $3$; two $2$s leave $15=3\cdot5$), the optimum is
\[
3+3+3+3+3+2+2=19,\qquad 3^5\cdot2^2=243\cdot4=\boxed{972}.
\]
(Equivalently $3^5\cdot4$.)
